% !TEX root = ../../../tree_notes.tex
% =====================================================================
%  Chapter 12 opener -- "Chapter at a Glance", three reading paths + the
%  ladder staircase, THE running example (the 10-node tree T, edges
%  1-2,1-3,2-4,2-5,3-6,3-7,5-8,7-9,9-10 -- promised six hand traces),
%  the "who owns what" comparison (metrics here vs Ch.1 binary recursions
%  vs Ch.16 query distances vs Ch.5 tours), roadmap, recognition guide,
%  constraint->tool table, notation.
%  Exports fig:c09-runtree and fig:c09-metricschart.
%  All tables below are hand-computed; every later section re-uses them.
% =====================================================================

\begin{chaptertheorybox}{Chapter at a Glance}
Chapters~7--11 put structure on \emph{values}: the BST orders them, the heap
keeps one extreme alive, the trie reads them symbol by symbol. This chapter asks a
different question -- forget what the nodes \emph{hold}; what does the tree's
\textbf{shape} itself owe you? Four families of answers, all computed by the same
two ingredients (a DFS with a parent parameter, and the farthest-node idea):

\term{measures} (\secref{c09-measures}): the \emph{depth} of a node, the
\emph{height} of a node, the \emph{size} of a subtree -- all three fall out of one
$\bigO{n}$ DFS, and every later chapter silently assumes you can produce them.
\term{diameter} (\secref{c09-diameter}): the longest path in the tree, by two
sweeps (\thmref{c09-twosweep}) or by DP joining the two deepest arms at each node
(\thmref{c09-dpjoin}) -- two proofs, one answer.
\term{centroid} (\secref{c09-centroid}): the node whose removal leaves no component
bigger than $n/2$ -- the tree's centre of mass, guaranteed to exist and to be one
or two adjacent nodes (\thmref{c09-centexist}).
\term{radius and centre} (\secref{c09-radius}): eccentricity, radius, the one-or-two
middle nodes of every diameter -- the ``where should the warehouse go'' answers.

\textbf{You will learn to}
\begin{itemize}[nosep]
  \item compute depth, height and subtree size for every node in one DFS on a
        \emph{general} tree, and rewrite it iteratively when the recursion depth
        is an input, not an accident; \hfill\lv{2}--\lv{3}
  \item prove and run both diameter algorithms -- two sweeps and
        join-the-two-arms DP -- and recover the actual path; \hfill\lv{3}--\lv{4}
  \item find a centroid in $\bigO{n}$ via subtree sizes and state exactly why the
        ``$n/2$'' threshold makes decomposition possible; \hfill\lv{3}--\lv{4}
  \item read eccentricity, radius and centre off a single pair of sweeps, and
        tell the centroid from the centre (they are \emph{different points}); \hfill\lv{4}
  \item combine everything on the Summit CF~1092F: the root that maximises a
        weighted distance sum, found by the reroot shift this chapter previews. \hfill\lv{5}
\end{itemize}

\textbf{Prerequisites.} DFS/BFS and the adjacency-list model from
Chapters~1--4 (\chref{c04} for rooting and the parent parameter); the bottom-up
recursion pattern of Chapter~2; a stack for the iterative rewrite. LCA-based
\emph{query} distance is deliberately \emph{not} assumed -- that is Chapter~16
(\chref{c11}); this chapter is one-shot computation.
\end{chaptertheorybox}

% ---------------- Reading paths + ladder position ----------------
\begin{chaptertheorybox}{How to Read This Chapter (three paths)}
Level badges $\lv1$--$\lv6$: $\lv2$ is a one-function measure you can write from
the definition, $\lv{3}$ is a solid LC-Medium you code straight from a section's
skeleton, $\lv{4}$ needs \emph{two} ideas composed (a sweep inside a prune, sizes
inside a centroid test), $\lv{5}$ is the Summit, where the invariant must hold
while the root moves.

\textbf{Fast path ($\lv2$--$\lv3$, ``I just need the tree's numbers'').} Read
\secref{c09-measures} and \secref{c09-diameter}; do the solved LC~742 and LC~1245
and the practices LC~2331, LC~1609. You leave able to cold-write depth/height/size
and both diameter algorithms -- which covers the entire ``tree metrics'' interview
band.

\textbf{Interview path ($\lv3$--$\lv4$).} Add \secref{c09-radius}: the
eccentricity-from-endpoints identity (\thmref{c09-eccend}) plus the centre rules
solve LC~310-style ``minimum height trees'' and the whole ``add a row / paint the
tree'' family; skim 12.3's centroid definition for the ``balance the split''
vocabulary interviews recycle.

\textbf{Full path (contest track).} Everything in order, then 12.3's centroid
\emph{existence} proof and the decomposition preview, then the Summit CF~1092F
and the mixed set -- at which point the forward rope is obvious: centroid
decomposition (Chapter~24, \chref{c18}), rerooting as theory (Chapter~15,
\chref{c20}), and distance \emph{queries} (Chapter~16, \chref{c11}).

The backward ropes: Chapter~4 (\chref{c04}) for rooting and LC~834's reroot shift
(quoted, not re-taught), Chapter~1 (\chref{c01}) for the binary height recursion
this chapter generalises.
\end{chaptertheorybox}

\begin{figure}[H]
\centering
\begin{tikzpicture}[diagram, x=9mm, y=7mm]
  \foreach \st/\lab in {1/lv1,2/lv2,3/lv3,4/lv4,5/lv5,6/lv6} {
    \fill[\lab!22, draw=\lab, line width=0.8pt] (\st-1,0) rectangle (\st,\st*0.52);
    \node[font=\scriptsize\bfseries, text=\lab] at (\st-0.5,\st*0.52-0.42) {L\st};}
  \node[annot, anchor=south west, font=\scriptsize]
    at (0.05,0.62) {depth, height, size -- one DFS};
  \node[annot, anchor=south west, font=\scriptsize]
    at (1.05,1.14) {12.1 iterative rewrite; LC 742};
  \node[annot, anchor=south west, font=\scriptsize]
    at (2.05,1.66) {12.2 two sweeps + DP join; LC 1245};
  \node[annot, anchor=south west, font=\scriptsize, text=lv5]
    at (3.05,2.18) {12.3 centroid proof; 12.4 centre rules};
  \node[annot, anchor=south west, font=\scriptsize, text=lv5]
    at (4.05,2.70) {12.4 CF 219D; reroot preview};
  \node[annot, anchor=south west, font=\scriptsize, text=lv6]
    at (5.05,3.22) {summit (CF 1092F)};
\end{tikzpicture}
\caption*{Chapter~12 on the six-step ladder: the three measures are $\lv2$--$\lv3$;
both diameter algorithms and the centre rules sit at $\lv3$--$\lv4$; the centroid
proof, CF~219D and every two-idea composition climb to $\lv4$--$\lv5$; the Summit
(CF~1092F, weighted distance sum by rerooting) is step~6 territory viewed from
step~5. Fast path $=$ steps 2--3; full path runs to 6.}
\end{figure}

% ---------------- The running example ----------------
\subsection*{The Running Example: One Ten-Node Tree, Six Angles}

\emph{One} unrooted tree does all the work. The nodes are $1,\dots,10$ and the
edges are
\[
  1\!-\!2,\; 1\!-\!3,\; 2\!-\!4,\; 2\!-\!5,\; 3\!-\!6,\; 3\!-\!7,\; 5\!-\!8,\; 7\!-\!9,\; 9\!-\!10 .
\]
Read it as a diagram first (\figref{c09-runtree}); the root used below is
\emph{a choice, not a fact} -- the tree has no preferred node until we pick one.

\begin{figure}[H]
\centering
\begin{tikzpicture}[diagram, scale=0.92, transform shape,
    tnode/.style={circle, draw=primaryblue, line width=0.9pt, fill=primaryblue!8,
                  minimum size=6.2mm, font=\small\bfseries, text=primaryblue!60!black}]
  \node[tnode] (n1) at (0,0) {$1$};
  \node[tnode] (n2) at (-3,-1.6) {$2$};
  \node[tnode] (n3) at ( 3,-1.6) {$3$};
  \node[tnode] (n4) at (-5.2,-3.2) {$4$};
  \node[tnode] (n5) at (-1.6,-3.2) {$5$};
  \node[tnode] (n6) at ( 1.6,-3.2) {$6$};
  \node[tnode] (n7) at ( 5.2,-3.2) {$7$};
  \node[tnode] (n8) at (-1.6,-4.8) {$8$};
  \node[tnode] (n9) at ( 5.2,-4.8) {$9$};
  \node[tnode] (n10) at ( 5.2,-6.4) {$10$};
  \draw[tedge] (n1)--(n2) (n1)--(n3) (n2)--(n4) (n2)--(n5)
               (n3)--(n6) (n3)--(n7) (n5)--(n8) (n7)--(n9) (n9)--(n10);
  \node[annot, font=\scriptsize, text=difforange!70!black, anchor=south east]
    at (-5.6,-3.45) {leaf};
  \node[annot, font=\scriptsize, text=difforange!70!black, anchor=north]
    at (5.2,-6.9) {leaf -- far end of the diameter $8\leftrightarrow10$ (7 edges)};
  \node[annot, font=\scriptsize, text=successgreen!50!black, anchor=west]
    at (1.2,0.35) {centroids $\{1,3\}$: cutting at either leaves pieces of $\le 5$ nodes};
\end{tikzpicture}
\caption{The running tree $T$ -- 10 nodes, 9 edges, unrooted. Rooted at~1 it has
height~4; its diameter is the path $8$--$5$--$2$--$1$--$3$--$7$--$9$--$10$ of
length~$7$; its centroids are the adjacent pair $1,3$ (the two-centroid case of
\thmref{c09-centexist}); its centre is the pair $\{1,3\}$. Every
section re-traces this same picture.}
\label{fig:c09-runtree}
\end{figure}

Six angles, six hand-computed tables -- each lives in the section named on its
row and never changes:

\begin{center}
\small
\renewcommand{\arraystretch}{1.25}
\begin{tabular}{@{}p{0.30\linewidth} p{0.44\linewidth} l@{}}
\toprule
\textbf{Angle on $T$} & \textbf{What comes out} & \textbf{Traced in} \\
\midrule
measures rooted at $1$ & depth $10{:}4$; height $1{:}4$; sizes $sz[2]=4$, $sz[3]=5$ & \secref{c09-measures} \\
two sweeps & from $4$: farthest is $10$ (dist $6$); from $10$: farthest is $8$ (dist $7$) & \secref{c09-diameter} \\
DP join & at $1$ the arms $3$ (via $2$) and $4$ (via $3$) join to $7$ & \secref{c09-diameter} \\
centroid cut & remove $1$: pieces of size $4$ and $5$, both $\le n/2$ --- and $3$ works too: two adjacent centroids & \secref{c09-centroid} \\
eccentricities & $\mathit{ecc}=[4,5,4,6,6,5,5,7,6,7]$, radius $4$, centre $\{1,3\}$ & \secref{c09-radius} \\
distance sums & $f(1)=20$ and the reroot shift gives $f(5)=28$, $f(10)=38$ & \secref{c09-radius}, Summit \\
\bottomrule
\end{tabular}
\end{center}

\begin{figure}[H]
\centering
\begin{tikzpicture}[diagram, scale=0.95, transform shape,
    mnode/.style={circle, draw=#1, line width=0.9pt, fill=#1!8, minimum size=5mm,
                  font=\scriptsize\bfseries, text=#1!60!black}]
  \node[mnode=primaryblue] (a) at (0,0)   {$d(v)$};
  \node[annot, font=\scriptsize, anchor=west] at (0.55,0) {$)=$ root-to-node edges \quad (12.1: depth)};
  \node[mnode=successgreen] (b) at (0,-0.85) {$h(v)$};
  \node[annot, font=\scriptsize, anchor=west] at (0.55,-0.85) {$=$ deepest node below $v$ \quad (12.1: height)};
  \node[mnode=difforange] (c) at (0,-1.7) {$sz(v)$};
  \node[annot, font=\scriptsize, anchor=west] at (0.55,-1.7) {$=$ nodes in $v$'s subtree \quad (12.1: size)};
  \node[mnode=warningred] (e) at (0,-2.55) {$D$};
  \node[annot, font=\scriptsize, anchor=west] at (0.55,-2.55) {$=$ longest path \quad (12.2: diameter $=$ two sweeps or DP)};
  \node[mnode=softpurple] (g) at (0,-3.4) {$C$};
  \node[annot, font=\scriptsize, anchor=west] at (0.55,-3.4) {$=$ max-piece $\le n/2$ \quad (12.3: centroid)};
  \node[mnode=primaryblue!60!black] (r) at (0,-4.25) {$R$};
  \node[annot, font=\scriptsize, anchor=west] at (0.55,-4.25) {$=\min_v\max(d(v,A),d(v,B))$ \quad (12.4: radius, $A,B$ diameter ends)};
\end{tikzpicture}
\caption*{The metric chart -- six numbers, two ingredients. Everything this
chapter computes is a $\mathit{max}$, $\mathit{min}$ or $\mathit{sum}$ over the
same DFS; the only new machinery is the \emph{two-sweep} trick behind $D$ and $R$.
}
\label{fig:c09-metricschart}
\end{figure}

% ---------------- Who owns what ----------------
\subsection*{Who Owns What (so nothing is learned twice)}

\begin{center}
\small
\renewcommand{\arraystretch}{1.25}
\begin{tabular}{@{}p{0.34\linewidth} p{0.30\linewidth} p{0.30\linewidth}@{}}
\toprule
\textbf{Idea} & \textbf{Owner} & \textbf{This chapter's share} \\
\midrule
height/size recursion on \emph{binary} trees & Ch.~1 (\thmref{c01-height-recursion}), Ch.~2 bottom-up & general trees, parent parameter, one DFS for \emph{all three} measures \\
rooting an undirected tree & Ch.~4 (\chref{c04}) & assumed; we add the metrics on top \\
distance per \emph{query} pair & Ch.~16 (\chref{c11}): LCA & one-shot \emph{all-node} sums instead \\
subtree as a contiguous range & Ch.~5 (\chref{c24}): Euler & sizes here are plain counts, no tour \\
rerooting theory & Ch.~15 (\chref{c20}) & one worked preview shift ($f(v)=f(u)+S-2sz[v]$) \\
centroid \emph{decomposition} & Ch.~24 (\chref{c18}) ``find \emph{one} centroid'' + why recursion is safe \\
\bottomrule
\end{tabular}
\end{center}

% ---------------- Recognition guide ----------------
\begin{recognitionbox}[Pattern card -- ``the problem names a shape number'']
\reckey{the statement asks for (or is blocked on) a \emph{number defined by the
tree's shape}: longest path, deepest node, biggest branch, ``best root'',
``minimum height'', ``farthest from everything''.}
\recuse{
  \item ``diameter / longest path / maximum distance'' -- both algorithms of
        \secref{c09-diameter}; the path-recovery variant reuses the same sweep
        with a parent map (\thmref{c09-twosweep});
  \item ``balance the tree / split into halves / safe recursion root'' -- the
        centroid of \secref{c09-centroid} (\defref{c09-centroid});
  \item ``place a facility / paint the tree / minimum height'' -- radius and
        centre, \thmref{c09-eccend} and \thmref{c09-centerrule};
  \item ``sum distances over all roots / choose the best capital'' -- the reroot
        shift previewed in \secref{c09-radius}; the Summit is its weighted
        completion.
}
\recnotuse{
  \item the edges \emph{carry} the answer (characters, bits) -- that is the trie,
        Chapter~11 (\chref{c08});
  \item you must answer distances for many \emph{live pairs of nodes} after
        preprocessing -- hand it to LCA, Chapter~16 (\chref{c11});
  \item the tree is \emph{binary} and the question is value-local (invert,
        validate, sum paths) -- Chapters~1--3, not metrics.
}
One-line test: \emph{does the answer depend on how far nodes are from each other,
not on what they store? Then it is a metric, and metrics live here.}
\end{recognitionbox}

% ---------------- Constraint -> tool ----------------
\subsection*{Reading the Constraints (which metric tool the limits want)}

\begin{center}
\small
\renewcommand{\arraystretch}{1.25}
\begin{tabular}{@{}p{0.30\linewidth} p{0.42\linewidth} l@{}}
\toprule
\textbf{Constraint signature} & \textbf{What it rules in/out} & \textbf{Tool} \\
\midrule
$n \le 10^{5}$--$10^{6}$ nodes, one-shot & any linear DFS family; \emph{not} all-pairs BFS & 12.1/12.2 sweeps \\
a \emph{chain} can be an input ($n \ge 2\cdot10^{5}$) & recursion depth is data-dependent -- budget the stack or go iterative & \coderef{c09-iter} \\
``longest path between any two nodes'' & diameter -- never try $\binom{n}{2}$ pairs & \thmref{c09-twosweep} \\
``remove a node, balance the rest'' / recursive on subtrees & centroid, not centre -- sizes, not distances & \thmref{c09-centfind} \\
``best root for a (weighted) distance sum'' & reroot shift over one initial DFS & Summit, \secref{c09-radius} \\
many distance \emph{queries} online & leave this chapter: LCA (Ch.~16) & \chref{c11} \\
\bottomrule
\end{tabular}
\end{center}

% ---------------- Notation ----------------
\subsection*{Notation Used Only Here}

$\depth(v)$ root-to-$v$ edges (once a root is fixed) $\mid$
$h(v)$ height of $v$'s subtree $\mid$ $sz(v)$ subtree node-count $\mid$
$\mathit{down}(v)$ longest downward \emph{arm} from $v$ $\mid$
$D$ the diameter $\mid$ $\mathit{ecc}(v)=\max_u d(v,u)$ $\mid$
radius $R=\min_v \mathit{ecc}(v)$ $\mid$ centroid: every component after the cut
has $\le \lfloor n/2\rfloor$ nodes $\mid$ $f(r)$ the cost/sum when the tree is
rooted at $r$; $S=\sum a_i$ the total weight (Summit).

Distances always count \emph{edges}; a single-node tree has diameter $0$. When the
tree is unrooted, ``rooted at $r$'' means the DFS view, and every claim below is
shown to be root-choice-independent where it must be.
